\section{Exam Required Extra Topics}

% Added from exam coverage
\subsection{Exam Formula Sheet Conventions}

\subsubsection{One-sided initial conditions}
\[
x(0^-),\quad \dot{x}(0^-),\quad y(0^-),\quad \dot{y}(0^-).
\]
Exam problems use unilateral Laplace transforms and left-limit initial conditions. These values enter the derivative rules before solving for \(Y(s)\).

% Added from exam coverage
\subsubsection{Determining whether an exam statement is false}
\[
\text{one false clause in an answer option} \quad \Rightarrow \quad \text{the option is false}.
\]
Several multiple-choice exams use answer options containing multiple claims. Every claim in an option must be true for the option to be true.

% Added from exam coverage
\subsection{Distributional Frequency Responses}

\subsubsection{Undamped oscillator frequency response}
\[
\mathcal{F}\{\sin \omega_0t\,u(t)\}
=\operatorname{PV}\frac{\omega_0}{\omega_0^2-\omega^2}
\frac{\pi}{2j}\left[\delta(\omega-\omega_0)-\delta(\omega+\omega_0)\right].
\]
Marginally stable systems can produce Fourier-domain expressions with principal-value and impulse terms. % TODO: verify from source

\subsubsection{Imaginary-axis poles and ordinary frequency response}
\[
H(s)=\frac{1}{s^2+\omega_0^2}
\quad \text{has poles at}\quad s=\pm j\omega_0.
\]
Substitution \(s=j\omega\) is not an ordinary finite frequency response at pole frequencies on the imaginary axis.
